% The two panels use the same risky policy and the same physical tree.
\begin{tikzpicture}[
  font=\figurefont\small,
  state/.style={draw=BookInk,fill=BookPaper,minimum width=12mm,minimum height=8mm},
  flow/.style={-{Stealth},BookInk,thick},
  backup/.style={-{Stealth},BookTeal,thick,dashed},
  every node/.style={align=center}]
  \node[anchor=west] at (0,3.4) {（a）联合概率上的静态CVaR};
  \node[state] (root) at (0.7,1.5) {出发};
  \node[state] (a) at (2.6,2.3) {A};
  \node[state] (b) at (2.6,0.5) {B};
  \node[state] (z) at (5.2,2.9) {0};
  \node[state] (bad) at (5.2,1.7) {20};
  \node[state] (ten) at (5.2,0.5) {10};
  \draw[flow] (root) -- node[above] {0.1} (a);
  \draw[flow] (root) -- node[below] {0.9} (b);
  \draw[flow] (a) -- node[above] {0.9} (z);
  \draw[flow] (a) -- node[below] {0.1} (bad);
  \draw[flow] (b) -- (ten);
  \node[anchor=west] at (5.95,2.9) {0.09};
  \node[anchor=west] at (5.95,1.7) {0.01};
  \node[anchor=west] at (5.95,0.5) {0.90};
  \node[draw=BookTeal,fill=BookTeal!6,minimum height=9mm]
    (static) at (3.5,-1) {尾部：$0.01\times20+0.09\times10$\\除以0.1，得到11};
  \draw[backup] (ten.south) -- (5.2,-0.3) -| (static.north);
  \begin{scope}[xshift=82mm]
    \node[anchor=west] at (0,3.4) {（b）逐阶段的嵌套CVaR};
    \node[state] (r2) at (0.7,1.5) {20};
    \node[state] (a2) at (2.6,2.3) {A：20};
    \node[state] (b2) at (2.6,0.5) {B：10};
    \node[state] (z2) at (5.2,2.9) {0};
    \node[state] (bad2) at (5.2,1.7) {20};
    \node[state] (ten2) at (5.2,0.5) {10};
    \draw[flow] (r2) -- node[above] {0.1} (a2);
    \draw[flow] (r2) -- node[below] {0.9} (b2);
    \draw[flow] (a2) -- node[above] {0.9} (z2);
    \draw[flow] (a2) -- node[below] {0.1} (bad2);
    \draw[flow] (b2) -- (ten2);
    \draw[backup] (bad2.south) -- (5.2,1.05) -| (a2.south);
    \draw[backup] (a2.north) -- (2.6,3.0) -| (r2.north);
    \node at (3.6,-1) {叶损失 $\longrightarrow$ A处尾部20\\分支风险 $\longrightarrow$ 根部尾部20};
  \end{scope}
\end{tikzpicture}
